\section{Validation, reproducibility, and formalization}
\label{sec:validation}

\subsection{What was checked}
The two new main arguments were independently reviewed within this analysis. The review checked the value-group multiplicity argument, simultaneous budget accounting, sharp binomial construction, integral monic division, interpolation denominators, extra variables, coordinate order lower bounds, and the edge cases $d=1$, $B=0$, and $M=1$.

Exact finite computations provide additional checks:
\begin{itemize}
\item \textbf{81 finite-field cases and 3,005 initial polynomials.} Multiplicities were computed by repeated exact polynomial division. A finite set-cover dynamic program found the minimum degree needed to cover all nonzero prime-field residues. Every result agreed with the proposed shell charge.
\item \textbf{12 shifted cyclotomic certificates.} Direct polynomial computations checked the Eisenstein initial relation used in the sharp construction.
\item \textbf{12 double-resultant certificates.} In examples $K=\Q_p(\pi)$ with $\pi^e=-p$, exact resultant polynomials and strict Newton-polygon inequalities certify that every conjugate obstruction value lies above its witness baseline.
\item \textbf{12 cutoff comparisons.} These include tame, wild, mixed, fractional-budget, and late-start examples, and check the strict budget crossing.
\item \textbf{15,624 exhaustive integer quadratic polynomials.} These test the finite-slope bad-choice bound with exact Gauss valuations.
\item \textbf{300 seeded packet checks.} Higher-degree examples over $p=2,3,5$ and four rational radii check simultaneous selection. The three-curve obstruction in Section~\ref{sec:applications} is also checked.
\end{itemize}
The finite checks are not proofs of the infinite theorems. Their role is to catch normalization, multiplicity, strict-threshold, and simultaneous-selection errors independently of the written derivation. The manuscript supplies the proofs.

\subsection{How to reproduce the artifacts}
From the archive root, the relevant commands are:
\begin{verbatim}
python3 code/torsion_bounds.py --p 2 --e 1 \
    --dimension 3 --rho 1 --budget 3
python3 code/check_zero_tolerance.py \
    --output checks/reproduced_zero_tolerance.json
python3 code/finite_slope_checks.py
pdflatex -interaction=nonstopmode -halt-on-error \
    grh_fixed_field_transfer.tex
pdflatex -interaction=nonstopmode -halt-on-error \
    grh_fixed_field_transfer.tex
\end{verbatim}
The bound calculator and slope checks use only Python's standard library. The zero-tolerance checks additionally require SymPy; its version is recorded in the JSON output. The slope script writes its JSON beside the script. Recorded outputs are preserved in \file{checks/}.

The calculator uses rational arithmetic and integer powers, including the ceiling of the logarithm via an integer comparison. Its output includes the selected starting level, final radius, scalar cost before and at the crossing, and every finite-slope range. It does not attempt to construct arithmetic interpolation data.

\subsection{Source and prior-work provenance}
The audit directory includes the pinned Git commit IDs, SHA-256 hashes of the source files directly studied, the SHA-256 values of all 31 incoming archives, and their text-member inventory. The source repository and the previous archives are not silently treated as one formal development. No new theorem in this report has been inserted into either remote repository.

The upstream checks described in its README concern the upstream formalization. The present report's new local theorems are proved on paper and supported by exact computations. They have not been implemented or kernel-checked in Lean in this session.

\subsection{A Lean formalization plan}
The scalar theorem requires four reusable lemmas:
\begin{enumerate}
\item Gauss valuation as a sum over polynomial roots, including the difference $G_\rho-G_0$.
\item Reduction at a prescribed radius and identification of its nonzero-root degree with the matching Newton-polygon multiplicity.
\item The divisibility of surviving exponent differences by $N_h/\gcd(e,N_h)$, followed by the characteristic-$p$ multiplicity bound.
\item A finite covering argument over $\F_p^\times$, and the explicit binomial sharpness construction.
\end{enumerate}
The multivariable theorem can then be formalized separately:
\begin{enumerate}
\item Integral division by a monic polynomial over a multivariable coefficient ring.
\item Gauss-valued Lagrange interpolation over its fraction field.
\item The unit binomial-coefficient estimate for $b_n-b_m$.
\item Counting bad slopes, induction on the number of variables, and the recurrence $u_{j+1}\le u_j^3$.
\item The final scalar selector, the lower order bound for $\zeta^{a_j}$, and truncation for formal series.
\end{enumerate}
The original \file{RootsOfUnity.lean} and \file{Lib/General/NumberTheory/GaussValuation.lean} provide relevant interfaces. The new claims should be stated with the coefficient field fixed before the order bound. A proof of the real inequality in Proposition~\ref{prop:rank} is elementary, but extracting the selected-twist version of the determinant theorem also requires preserving the existing height and conjugate-integrality hypotheses.

This plan deliberately distinguishes mathematical readiness from completed formal verification. The archive contains no placeholder Lean theorem presented as a checked proof.

